\fsection[\textcolor{waveBlue2}{Ejercicio 3}]{\boldit{\textcolor{waveBlue2}{Actividades}}{
Señala en tu cuaderno si las siguientes afirmaciones son verdaderas o falsas . Justifica tu respuesta
}}

\begin{solution}

\textbf{La economía verde se preocupa solo de los objetivos medioambientales $\to$ \textcolor{samuraiRed}{\textbf{Falso}}.}
\begin{itemize}
    \item La economía verde integra dimensiones ambientales, económicas y sociales.
\end{itemize}

\textbf{El principio de justicia hace referencia a que los gobiernos deben basarse en evidencia científica $\to$ \textcolor{samuraiRed}{\textbf{Falso}}.}
\begin{itemize}
    \item Ese enfoque corresponde a la buena gobernanza; la justicia se refiere a la equidad en la distribución de recursos y oportunidades.
\end{itemize}

\textbf{El principio de los límites del planeta implica que las actividades económicas deben hacer frente a los grandes retos medioambientales $\to$ \textcolor{springGreen}{\textbf{Verdadero}}.}
\begin{itemize}
    \item La actividad económica debe respetar los límites ecológicos del planeta.
\end{itemize}

\textbf{El principio de bienestar pone de manifiesto que debemos reducir el consumo de bienes y servicios $\to$ \textcolor{samuraiRed}{\textbf{Falso}}.}
\begin{itemize}
    \item El bienestar se centra en mejorar la calidad de vida de forma sostenible, no únicamente en reducir consumo.
\end{itemize}

\textbf{El principio de buena gobernanza requiere que las instituciones sean transparentes $\to$ \textcolor{springGreen}{\textbf{Verdadero}}.}
\begin{itemize}
    \item Implica transparencia, participación y toma de decisiones basada en criterios objetivos.
\end{itemize}
\end{solution}
